\documentclass[11pt, a4paper]{article}
\usepackage[margin=2.5cm, headheight=14pt]{geometry}
\usepackage{titlesec}
\usepackage{booktabs}
\usepackage{longtable}
\usepackage{array}
\usepackage{enumitem}
\usepackage{xcolor}
\usepackage{fancyhdr}
\usepackage{hyperref}
\usepackage{listings}
\usepackage[T1]{fontenc}
\usepackage{parskip}
\usepackage{mdframed}
\usepackage{tabularx}
\usepackage{multirow}

% Colors
\definecolor{accent}{RGB}{30, 80, 162}
\definecolor{warn}{RGB}{180, 60, 40}
\definecolor{codebg}{RGB}{245, 245, 245}
\definecolor{darkgray}{RGB}{60, 60, 60}
\definecolor{green}{RGB}{30, 120, 60}

% Section formatting
\titleformat{\section}{\large\bfseries\color{accent}}{}{0em}{}[\titlerule]
\titleformat{\subsection}{\normalsize\bfseries\color{darkgray}}{}{0em}{}
\titleformat{\subsubsection}{\normalsize\itshape}{}{0em}{}

% Code style
\lstset{
  backgroundcolor=\color{codebg},
  basicstyle=\ttfamily\scriptsize,
  breaklines=true,
  breakatwhitespace=false,
  frame=single,
  framesep=4pt,
  rulecolor=\color{gray!40},
  xleftmargin=8pt,
  xrightmargin=8pt,
  columns=flexible,
}

% Header/Footer
\pagestyle{fancy}
\fancyhf{}
\rhead{\textcolor{gray}{\footnotesize Project A1 --- Freight Cost Intelligence}}
\lhead{\textcolor{gray}{\footnotesize Sameer --- Portfolio Roadmap}}
\rfoot{\textcolor{gray}{\footnotesize \thepage}}
\lfoot{\textcolor{gray}{\footnotesize Apr--Jul 2026}}
\renewcommand{\headrulewidth}{0.4pt}
\renewcommand{\footrulewidth}{0.4pt}

% Framed boxes
\newmdenv[
  linecolor=accent,
  linewidth=1.5pt,
  topline=true, bottomline=true, leftline=true, rightline=false,
  backgroundcolor=accent!5,
  innerleftmargin=10pt, innerrightmargin=10pt,
  innertopmargin=8pt, innerbottommargin=8pt,
]{infobox}

\newmdenv[
  linecolor=warn,
  linewidth=1.5pt,
  topline=true, bottomline=true, leftline=true, rightline=false,
  backgroundcolor=warn!5,
  innerleftmargin=10pt, innerrightmargin=10pt,
  innertopmargin=8pt, innerbottommargin=8pt,
]{warnbox}

\newmdenv[
  linecolor=green,
  linewidth=1.5pt,
  topline=true, bottomline=true, leftline=true, rightline=false,
  backgroundcolor=green!5,
  innerleftmargin=10pt, innerrightmargin=10pt,
  innertopmargin=8pt, innerbottommargin=8pt,
]{featurebox}

\hypersetup{colorlinks=true, linkcolor=accent, urlcolor=accent}

\begin{document}

% ── Title Page ──────────────────────────────────────────────────────────────
\begin{titlepage}
  \centering
  \vspace*{2cm}
  {\huge\bfseries\color{accent} Project A1\par}
  \vspace{0.5cm}
  {\LARGE Freight Cost Intelligence\par}
  \vspace{0.3cm}
  {\large Health Commodity Freight Cost Prediction Enriched with UNCTAD Bilateral Rate Benchmarks\par}
  \vspace{0.2cm}
  {\normalsize USAID Transaction Data + UNCTAD-World Bank Global Transport Costs Dataset + Unified RAG + What-If Chat\par}
  \vspace{2cm}
  \rule{0.6\textwidth}{0.4pt}\\[1em]
  {\large Sameer \textbar{} Habib University, Karachi\par}
  {\large Portfolio Block 1 --- April to July 2026\par}
  \vspace{1cm}
  {\normalsize Part of a 4-project arc targeting supply chain data science roles in Pakistan\\
  and later RWTH Aachen MSc Data Science application (Winter 2028/29)\par}
  \vspace{2cm}
  \rule{0.6\textwidth}{0.4pt}\\[1em]
  {\small\color{gray} v4 --- Pivoted from two-dataset multimodal to USAID + UNCTAD enrichment after Brunel EDA gate failure (single-date snapshot, 99.5\% AIR imbalance, 24\% unratable orders)}
\end{titlepage}

\tableofcontents
\newpage

% ── 1. Why This Exists ───────────────────────────────────────────────────────
\section{Why This Project Exists}

\begin{infobox}
\textbf{One-line goal:} Build and deploy a health commodity freight cost prediction model on USAID shipment data, enriched with UNCTAD bilateral freight rate benchmarks as external features, with generative SHAP-grounded explanation, cited trade-doc context, and a what-if chat --- proving end-to-end SC-DS + LLM integration capability for Pakistan junior analyst roles and RWTH/Werkstudent applications.
\end{infobox}

This is not a Kaggle competition. The goal is a shipped, deployed artifact that a recruiter can click, a hiring manager can interrogate, and an RWTH admissions committee reads as a concrete portfolio piece.

\textbf{Why supply chain?} Supply chain was chosen over generic DS and fintech as the primary career vertical. It pairs domain specificity with genuine demand in Pakistan's logistics, e-commerce, and FMCG sectors. Junior analyst roles at companies like Daraz, Nestl\'{e}, and 3PLs explicitly look for cost modelling and demand forecasting skills.

\textbf{Why generative explanation + chat?} A plain prediction number has zero interview story. A prediction with a cited plain-English explanation --- grounded in the model's SHAP values and retrieved trade documents --- is memorable, differentiated, and demonstrates LLM integration, which is increasingly expected in SC-DS roles. The what-if chat upgrades the tool from demo to decision-support, which is the actual job.

\textbf{Why USAID + UNCTAD enrichment, not two separate models?} The v3 plan included Brunel Supply Chain Logistics as a second dataset. Brunel EDA (Week 2) revealed fatal structural problems: all 9,215 orders share a single date (no time-based split possible), 99.5\% of ratable orders are AIR (GROUND sub-model has 34 rows), and 24\% of orders are unratable due to a discontinued carrier. Brunel is a linear-programming benchmark snapshot, not a freight cost time series. It was dropped at the Week 2 decision gate.

The replacement enrichment strategy uses the \textbf{UNCTAD-World Bank Global Transport Costs Dataset for International Trade (GTCDIT)}, a free, authoritative dataset covering bilateral freight rates (air, sea, rail, road) by country pair and commodity group, 2016--2021. This is used not as a second training dataset but as a \textbf{lookup feature table}: for each USAID shipment, we join the UNCTAD ad-valorem freight rate benchmark for the relevant origin country $\rightarrow$ destination country $\rightarrow$ mode combination. This gives the model an external structural cost signal --- the kind of benchmark a real freight analyst would consult. The framing in the README is accurate and defensible: ``We use USAID transaction data enriched with UNCTAD bilateral freight rate benchmarks as a lane-level cost signal.''

\begin{warnbox}
\textbf{Temporal mismatch --- know this for interviews.} USAID spans approximately 2006--2015; UNCTAD covers 2016--2021. There is no year overlap. UNCTAD is therefore used as a \textbf{structural benchmark}, not a time-matched join. The assumption: bilateral freight rate intensities for a given lane and mode are sufficiently stable structural signals even when drawn from a different period. This is standard practice in gravity-model freight research. In the README and in interviews, state this explicitly. Do not pretend the join is temporal.
\end{warnbox}

\textbf{How this fits the larger plan:}
\begin{itemize}[noitemsep]
  \item Phase 1 (now): Build portfolio, GitHub, LinkedIn presence.
  \item Phase 2 (4th year + 1yr post-grad): Land SC data analyst/DS role in Pakistan (80--150k PKR).
  \item Phase 3 (mid-2027): GRE, IELTS, APS certificate, blocked account (EUR 11,208).
  \item Phase 4: RWTH Aachen MSc Data Science (Winter 2028/29) $\rightarrow$ Werkstudent $\rightarrow$ found company.
\end{itemize}

Seeds the final-year project (SC data + RAG/ML + deployed tool = triple asset: portfolio, thesis demo, interview piece).

% ── 2. Project Specification ─────────────────────────────────────────────────
\section{Project Specification}

\subsection{Problem Statement}
Build a freight cost prediction model on USAID health commodity shipment data (Air/Truck/Ocean, 10k+ records), enriched with UNCTAD bilateral freight rate benchmarks as engineered features. The model returns a predicted cost plus uncertainty range. A generative explanation layer reads SHAP values and retrieves relevant trade-document chunks to produce a plain-English summary with feature breakdown and cited context. A what-if chat lets the user ask follow-ups including scenario modifications (``what if I switch mode?'', ``what if destination is Nigeria instead?'') that rerun the model and return a new prediction with explanation.

\subsection{App Features (Final State)}

\begin{featurebox}
\textbf{Tab 1 --- Predict and Explain}
\begin{enumerate}[noitemsep]
  \item User fills shipment form: origin hub, destination country, mode (Air/Truck/Ocean), weight, commodity, INCO term.
  \item Model returns predicted freight cost + uncertainty range.
  \item UNCTAD benchmark for the lane is shown alongside: ``UNCTAD ad-valorem rate for [origin] $\rightarrow$ [destination] via [mode]: X\%''.
  \item LLM generates plain-English summary of why this cost was predicted.
  \item Feature importance breakdown table shown (top 5--7 SHAP drivers), including UNCTAD feature if it appears.
  \item Cited trade-doc context shown (e.g. relevant Incoterms clause, HS code note).
\end{enumerate}
\end{featurebox}

\begin{featurebox}
\textbf{Tab 2 --- What-If Chat}
\begin{enumerate}[noitemsep]
  \item Chat pre-seeded with current shipment context from Tab 1.
  \item User asks alternatives: ``What if I use sea instead of air?'' ``What if I ship to Uganda instead of Kenya?''
  \item System parses intent $\rightarrow$ modifies input dict (including re-looking up UNCTAD benchmark for new lane) $\rightarrow$ reruns model $\rightarrow$ returns new prediction.
  \item Response includes: new cost, delta vs original, plain-English explanation, cited trade docs.
  \item General trade questions (Incoterms, HS codes, glossary) answered via unified RAG with citations.
\end{enumerate}
\end{featurebox}

\subsection{Datasets}

\begin{tabularx}{\textwidth}{lXl}
\toprule
\textbf{Dataset} & \textbf{Description} & \textbf{Source} \\
\midrule
USAID Supply Chain Shipment Pricing & 10k+ real shipment records: origin (hub), destination country, mode (Air dominant, Truck, some Ocean, N/A), weight, freight cost, INCO term, commodity. Domain: HIV/ARV health commodities to PEPFAR countries. Time range: approx. 2006--2015. Used as primary training dataset. & Kaggle (apoorvwatsky / divyeshardeshana) \\
\addlinespace
UNCTAD-World Bank GTCDIT & Global Transport Costs Dataset for International Trade. Bilateral freight rate benchmarks (ad-valorem \%, \$/kg) by origin country, destination country, commodity group (HS 2017 4-digit), and mode (air, sea, rail, road). Covers 2016--2021, 170+ economies, 92\% of global merchandise trade. Used as a lookup feature table joined onto USAID records by country-pair + mode. Not a training dataset --- a structural benchmark enrichment. & UNCTADstat Data Centre: \url{https://unctadstat.unctad.org/datacentre/dataviewer/US.TransportCosts} (free, bulk download available) \\
\addlinespace
RAG Corpus & Incoterms 2020 explainers, HS code chapter descriptions, shipping glossaries. Mode-agnostic trade knowledge. & ICC, WCO, public freight sites \\
\bottomrule
\end{tabularx}

\begin{warnbox}
\textbf{Do not commit raw data to git.} Add \texttt{data/raw/} and \texttt{data/processed/} to \texttt{.gitignore}. Corpus PDFs in \texttt{data/corpus/} can be committed if under 10MB total.
\end{warnbox}

\begin{infobox}
\textbf{On the UNCTAD enrichment design:} UNCTAD GTCDIT is a lookup table, not a training dataset. For each USAID shipment row, we look up the UNCTAD ad-valorem freight rate for (origin\_country, destination\_country, mode) and add it as a feature column (\texttt{unctad\_adval\_rate}). If no match exists for a specific HS code, aggregate to commodity group or fall back to mode-level average for that country pair. Missing coverage $\rightarrow$ NaN $\rightarrow$ LightGBM handles natively. This produces one model, one training dataset, one deployment --- significantly simpler architecture than the two-model plan.
\end{infobox}

\subsection{Stack}

\begin{tabularx}{\textwidth}{lX}
\toprule
\textbf{Layer} & \textbf{Tools} \\
\midrule
Data wrangling & \texttt{pandas}, \texttt{numpy} \\
Modelling & \texttt{scikit-learn} (pipelines, baseline), \texttt{xgboost}, \texttt{lightgbm} \\
Explainability & \texttt{shap} --- SHAP values for feature importance, feeds LLM explanation prompt \\
Visualization & \texttt{matplotlib}, \texttt{seaborn} \\
Embeddings & \texttt{sentence-transformers} (\texttt{all-MiniLM-L6-v2}) --- CPU-safe \\
Vector store & \texttt{faiss-cpu} (primary) \textbar{} \texttt{chromadb} (alt) \\
PDF parsing & \texttt{pypdf} (primary) \textbar{} \texttt{pdfplumber} (alt) \\
LLM client & \texttt{groq} (free, primary) \textbar{} \texttt{openai} (fallback via env var) \\
App & \texttt{streamlit} \\
Deployment & Hugging Face Spaces (free) \\
Secrets & \texttt{python-dotenv} local, HF Secrets in prod \\
Tests & \texttt{pytest} \\
Linting & \texttt{ruff} \\
\bottomrule
\end{tabularx}

\textbf{Key addition: \texttt{shap}.} SHAP values are the bridge between the ML model and the LLM explanation. The top SHAP contributors for a given prediction get injected into the LLM prompt, grounding the explanation in actual model reasoning rather than hallucination. If \texttt{unctad\_adval\_rate} has high SHAP importance, the explanation can cite it as ``the UNCTAD benchmark rate for this lane was X\%, signalling a high-cost corridor.''

\textbf{Hardware:} 16GB RAM, Intel HD only. All choices are CPU-safe. Do not use CUDA-dependent models.

% ── 3. Repository Structure ──────────────────────────────────────────────────
\section{Repository Structure}

\begin{lstlisting}
freight-cost-rag/
├── README.md
├── requirements.txt
├── .gitignore
├── .env.example
├── data/
│   ├── raw/
│   │   ├── usaid/          # SCMS_Delivery_History CSV
│   │   └── unctad/         # GTCDIT bulk download CSV(s)
│   ├── processed/          # gitignored
│   └── corpus/             # RAG PDFs -- commit if small
├── notebooks/
│   ├── 01_eda_usaid.ipynb
│   ├── 02_eda_unctad.ipynb        # schema tour, coverage audit, join plan
│   ├── 03_feature_engineering.ipynb
│   ├── 04_baseline.ipynb
│   ├── 05_tree_models.ipynb
│   ├── 06_unctad_feature_analysis.ipynb  # SHAP analysis of UNCTAD feature
│   └── 07_rag_prototype.ipynb
├── src/
│   ├── __init__.py
│   ├── config.py
│   ├── data/
│   │   ├── __init__.py
│   │   ├── load_usaid.py
│   │   └── load_unctad.py    # loads GTCDIT, builds lookup dict
│   ├── features/
│   │   ├── __init__.py
│   │   └── features.py       # USAID features + UNCTAD enrichment join
│   ├── models/
│   │   ├── __init__.py
│   │   ├── train.py
│   │   └── predict.py
│   ├── explainer.py          # SHAP + LLM explanation
│   └── rag/
│       ├── __init__.py
│       ├── ingest.py
│       ├── retriever.py
│       ├── generator.py
│       └── whatif.py
├── tests/
│   ├── test_features.py
│   ├── test_unctad_join.py   # verify coverage, fallback logic
│   ├── test_explainer.py
│   └── test_rag.py
├── app/
│   └── streamlit_app.py
├── models/                   # gitignored
│   └── usaid_model.pkl
└── scripts/
    └── train.py
\end{lstlisting}

\subsection{Copy-Paste Files}

\subsubsection{.gitignore}
\begin{lstlisting}
__pycache__/
*.pyc
.env
.venv/
venv/
data/raw/
data/processed/
models/
*.ipynb_checkpoints/
.DS_Store
.vscode/
.idea/
*.log
.streamlit/secrets.toml
\end{lstlisting}

\subsubsection{requirements.txt}
\begin{lstlisting}
pandas>=2.0
numpy
scikit-learn>=1.3
xgboost
lightgbm
shap
matplotlib
seaborn
jupyter
sentence-transformers
faiss-cpu
pypdf
groq
streamlit
python-dotenv
pytest
ruff
\end{lstlisting}

\subsubsection{.env.example}
\begin{lstlisting}
GROQ_API_KEY=your_key_here
OPENAI_API_KEY=optional_fallback
LLM_PROVIDER=groq
\end{lstlisting}

% ── 4. 13-Week Roadmap ───────────────────────────────────────────────────────
\section{13-Week Roadmap}

\textbf{Schedule:} Weekend-heavy. Saturday/Sunday = 4--5hr deep work. Weekdays = 30--45min touches. Do not fight this rhythm.

\begin{infobox}
\textbf{Status at time of writing (v4):} Week 1 USAID EDA complete. Brunel dropped at Week 2 decision gate. Plan restarts at Week 2 with UNCTAD EDA. USAID findings already documented: time-based 80/20 split at Dec 2013; weight = commodity weight (unit $\times$ qty); 5 outliers need 99th-pct capping; freight costs doubled over dataset period; log-log weight vs cost shows power-law relationship; CIF/DAP/DDU grouped as ``Other''.
\end{infobox}

\subsection{Phase 1 --- UNCTAD EDA + USAID Feature Engineering + Baseline (Weeks 2--4)}

\subsubsection*{Week 2 --- UNCTAD Pull and EDA (\textasciitilde{}7 hrs)}

\textbf{Download:}
\begin{itemize}[noitemsep]
  \item Go to \url{https://unctadstat.unctad.org/datacentre/dataviewer/US.TransportCosts}
  \item Select: all origins, all destinations, mode = air + sea (the two modes relevant to USAID), years = all available (2016--2021), indicator = ad-valorem transport costs (\%) and per-unit transport costs (\$/kg).
  \item Download bulk CSV $\rightarrow$ \texttt{data/raw/unctad/}.
\end{itemize}

\textbf{Weekday touches --- \texttt{src/data/load\_unctad.py}:}
\begin{itemize}[noitemsep]
  \item Load GTCDIT CSV. Understand columns: origin (M49 code), destination (M49 code), commodity (HS 2017 4-digit), mode, year, ad-valorem rate, \$/kg rate.
  \item Map M49 country codes to ISO-3 or country names to match USAID destination country column.
  \item Build a lookup dict: \texttt{(origin\_country, dest\_country, mode) $\rightarrow$ median ad-valorem rate across years}. Median across 2016--2021 gives a stable structural benchmark.
  \item Check: what share of USAID origin/destination/mode combinations have UNCTAD coverage?
\end{itemize}

\textbf{Weekend block --- \texttt{02\_eda\_unctad.ipynb}:}
\begin{itemize}[noitemsep]
  \item Coverage audit: for USAID rows, what \% get a direct UNCTAD match? What \% require country-level aggregation fallback? What \% remain NaN?
  \item Distribution of UNCTAD ad-valorem rates for USAID lanes. Are the rates plausible (air should be higher than sea)?
  \item Correlation: does UNCTAD benchmark correlate with USAID actual freight cost at the lane level?
  \item ``Findings'' markdown cell: honest read on UNCTAD coverage and signal quality.
\end{itemize}

\begin{warnbox}
\textbf{Week 2 decision gate for UNCTAD:} If coverage is below 30\% of USAID rows even after fallback, or if UNCTAD ad-valorem rates show zero correlation with USAID freight costs, drop the UNCTAD feature. Reframe to USAID-only. Do not sunk-cost a weak external feature. Post findings before proceeding.
\end{warnbox}

\subsubsection*{Week 3 --- Feature Engineering (\textasciitilde{}7 hrs)}

\textbf{Weekday touches:}
\begin{itemize}[noitemsep]
  \item \texttt{src/data/load\_usaid.py} --- clean loader (outlier capping, mode grouping, log target).
  \item \texttt{src/features/features.py} --- route encoding, weight bucketing, mode handling, INCO term encoding, year feature, UNCTAD join logic with fallback hierarchy.
  \item \texttt{tests/test\_features.py} --- sanity checks.
  \item \texttt{tests/test\_unctad\_join.py} --- verify fallback logic: HS4 $\rightarrow$ HS2 $\rightarrow$ mode-level average $\rightarrow$ NaN.
\end{itemize}

\textbf{Weekend block --- \texttt{03\_feature\_engineering.ipynb}:}
\begin{itemize}[noitemsep]
  \item Final feature set with UNCTAD column included.
  \item Missing value audit post-join. How many NaNs in \texttt{unctad\_adval\_rate}?
  \item Time-based split: 80\% train (pre-Dec 2013) / 20\% test (post-Dec 2013).
  \item README update: problem, data, enrichment approach.
\end{itemize}

\subsubsection*{Week 4 --- Baseline + Tree Models (\textasciitilde{}7 hrs)}

\textbf{\texttt{04\_baseline.ipynb}:}
\begin{itemize}[noitemsep]
  \item Log-transform target. LinearRegression baseline (with and without UNCTAD feature). RMSE, MAE, R\textsuperscript{2}.
\end{itemize}

\textbf{\texttt{05\_tree\_models.ipynb}:}
\begin{itemize}[noitemsep]
  \item XGBoost + LightGBM. LightGBM native categorical handling for country, INCO term, commodity.
  \item LightGBM handles NaN in \texttt{unctad\_adval\_rate} natively --- no imputation needed.
  \item Compare vs linear baseline.
  \item One round \texttt{RandomizedSearchCV} on LightGBM. No rabbit-holing.
  \item Pickle $\rightarrow$ \texttt{models/usaid\_model.pkl}.
  \item Install \texttt{shap}. Compute SHAP on 5 sample predictions. Check whether \texttt{unctad\_adval\_rate} appears in top-10 features.
\end{itemize}

\subsection{Phase 2 --- UNCTAD Feature Analysis + Buffer (Weeks 5--6)}

\subsubsection*{Week 5 --- UNCTAD Feature Analysis (\textasciitilde{}7 hrs)}

\textbf{\texttt{06\_unctad\_feature\_analysis.ipynb}:}
\begin{itemize}[noitemsep]
  \item SHAP importance of \texttt{unctad\_adval\_rate} vs other features.
  \item Partial dependence plot: as UNCTAD benchmark increases, does predicted cost increase? (Sanity check.)
  \item Residuals by UNCTAD coverage: does the model perform better on rows with direct UNCTAD match vs fallback vs NaN?
  \item Honest writeup: what does the UNCTAD feature contribute, and what are its limitations (temporal mismatch, aggregation)?
\end{itemize}

\subsubsection*{Week 6 --- Buffer (\textasciitilde{}7 hrs)}

\begin{warnbox}
\textbf{Non-negotiable buffer.} Absorbs whatever broke in Weeks 2--5. If genuinely ahead: refactor \texttt{src/models/} into clean \texttt{train()} + \texttt{predict(inputs)} functions. Write \texttt{scripts/train.py} as CLI entry point.
\end{warnbox}

\subsection{Phase 3 --- RAG Layer (Weeks 7--9)}

\subsubsection*{Week 7 --- Corpus and Ingestion}

Build \texttt{data/corpus/}:
\begin{itemize}[noitemsep]
  \item Incoterms 2020 explainers: ICC public summaries + freight forwarder PDFs (5--10 docs).
  \item HS code chapter descriptions: WCO public or CSV-to-text.
  \item Shipping glossary: logistics firm websites, public PDFs.
\end{itemize}

\texttt{src/rag/ingest.py}: parse PDFs $\rightarrow$ chunk (500 tokens, 50 overlap) $\rightarrow$ embed (\texttt{all-MiniLM-L6-v2}) $\rightarrow$ FAISS index $\rightarrow$ save to disk.

Alt if PDFs chaotic: convert corpus to \texttt{.txt} manually once, commit, parse as plain text.

\subsubsection*{Week 8 --- Retriever and Generator}

\texttt{src/rag/retriever.py}: load FAISS, similarity search, return top-k chunks with source metadata.

\texttt{src/rag/generator.py}: Groq client, prompt template, answer + citations.

\begin{lstlisting}[language=Python]
# src/config.py
import os

LLM_PROVIDER = os.getenv("LLM_PROVIDER", "groq")

def get_llm_client():
    if LLM_PROVIDER == "groq":
        from groq import Groq
        return Groq(api_key=os.getenv("GROQ_API_KEY"))
    elif LLM_PROVIDER == "openai":
        from openai import OpenAI
        return OpenAI(api_key=os.getenv("OPENAI_API_KEY"))
\end{lstlisting}

\texttt{07\_rag\_prototype.ipynb}: 10 test queries end-to-end. Log quality. Note failures.

\subsubsection*{Week 9 --- Generative Explanation (SHAP $\rightarrow$ LLM)}

Implement \texttt{src/explainer.py}. This is where ML and RAG connect.

\textbf{Pipeline:}
\begin{enumerate}
  \item User submits form.
  \item Model predicts cost. SHAP computes top-5 feature contributions for this prediction.
  \item Retriever pulls relevant trade doc chunks.
  \item LLM prompt assembled: shipment inputs + SHAP table + retrieved chunks + UNCTAD benchmark for this lane.
  \item LLM returns plain-English summary + cites trade docs.
\end{enumerate}

\textbf{Prompt template:}
\begin{lstlisting}
SYSTEM: You are a freight cost analyst specialising in health commodity
supply chains. Explain shipment cost predictions clearly. Always cite
the trade documents provided. Never invent rules not present in context.

USER:
Shipment: {origin} to {destination}, mode: {mode}, weight: {weight}kg,
commodity: {commodity}, INCO term: {inco_term}.

Predicted cost: ${predicted_cost} (range: ${low} -- ${high})

UNCTAD bilateral freight benchmark for {origin} -> {destination} via {mode}:
Ad-valorem rate: {unctad_adval_rate}% (2016-2021 median, structural benchmark)

Top cost drivers (SHAP values, higher = more impact on cost):
{shap_table}

Relevant trade document excerpts:
{retrieved_chunks}

Write:
1. A 3-4 sentence plain-English summary of why this cost was predicted.
2. A bullet list of the top cost drivers with a one-line explanation each.
3. One relevant trade rule or definition from the documents above
   (cite the source).
\end{lstlisting}

Add query caching (\texttt{functools.lru\_cache} or dict keyed on input hash).

\subsection{Phase 4 --- What-If Chat and App (Weeks 10--12)}

\subsubsection*{Week 10 --- What-If Chat Logic}

Implement \texttt{src/rag/whatif.py}.

\textbf{Architecture:}
\begin{enumerate}
  \item \textbf{Intent parsing:} LLM call \#1 --- extract structured modification intent from user message + current shipment context.
  \item \textbf{Input modification:} apply changes; re-lookup UNCTAD benchmark for new lane if origin/destination/mode changed.
  \item \textbf{Reprediction:} run modified inputs through model. New cost + SHAP.
  \item \textbf{RAG retrieval:} retrieve trade doc chunks relevant to the modification.
  \item \textbf{Response generation:} LLM call \#2 --- new prediction, delta, explanation, citations.
\end{enumerate}

\textbf{Intent parsing prompt:}
\begin{lstlisting}
SYSTEM: Extract shipment modification intent as JSON. Only output JSON.
No preamble.

USER:
Current shipment: {shipment_dict_as_json}
User message: "{user_message}"

Output JSON with keys:
{
  "modification_type": "mode_change | route_change | weight_change | other",
  "changes": {"field": "new_value"},
  "is_what_if": true | false,
  "rag_only": true | false
}
\end{lstlisting}

\textbf{Stateful chat in Streamlit:} use \texttt{st.session\_state} to persist shipment dict and chat history across reruns.

\subsubsection*{Week 11 --- Streamlit App (Local)}

\texttt{app/streamlit\_app.py} --- two tabs:

\textbf{Tab 1 --- Predict and Explain:}
\begin{itemize}[noitemsep]
  \item Shipment form: origin hub, destination country, mode, weight, commodity, INCO term.
  \item Predict $\rightarrow$ cost + uncertainty.
  \item Explanation: LLM summary + SHAP table + UNCTAD benchmark callout + cited trade context.
  \item ``Open chat about this shipment'' $\rightarrow$ Tab 2, pre-seeded.
  \item \texttt{@st.cache\_resource} for model + FAISS index + UNCTAD lookup dict.
\end{itemize}

\textbf{Tab 2 --- What-If Chat:}
\begin{itemize}[noitemsep]
  \item Shows current shipment summary at top.
  \item Chat input, history via \texttt{st.chat\_message}.
  \item Example prompts: ``What if I use sea freight?'', ``What if destination is Uganda?'', ``What does CIF mean?''
  \item ``Reset shipment'' button clears session state.
\end{itemize}

\subsubsection*{Week 12 --- Deploy and Polish}

\begin{enumerate}[noitemsep]
  \item Create HF Space (Streamlit), link to GitHub.
  \item Add \texttt{GROQ\_API\_KEY} to HF Secrets.
  \item Cold start test. Document wait in README.
  \item Verify model pickle + FAISS + UNCTAD lookup dict load on HF CPU.
  \item Example shipments (one-click ``Try this'').
  \item Friendly error messages for rate limits.
  \item ``How it works'' collapsible: diagram showing USAID data + UNCTAD enrichment + RAG + LLM.
\end{enumerate}

\subsection{Phase 5 --- Ship (Week 13)}

\subsubsection*{Week 13 --- Documentation + Interview Prep}

README 30-second skim test:
\begin{enumerate}[noitemsep]
  \item Problem: health commodity freight cost prediction.
  \item Data: USAID transaction records enriched with UNCTAD bilateral freight rate benchmarks.
  \item Approach: LightGBM with UNCTAD enrichment, RAG-grounded generative explanation, what-if chat.
  \item Results: RMSE/MAE, baseline vs LightGBM, UNCTAD feature SHAP importance.
  \item Demo: link + screenshots.
  \item Next steps: Olist last-mile dataset as A2, commercial freight data integration.
\end{enumerate}

\textbf{Interview Q\&A (14 samples):}
\begin{itemize}[noitemsep]
  \item Why USAID + UNCTAD enrichment instead of two separate ML datasets?
  \item Why was Brunel dropped? What did the EDA reveal?
  \item What is the temporal mismatch between USAID and UNCTAD, and how do you handle it?
  \item Why log-transform the target?
  \item How does SHAP ground the LLM explanation?
  \item Why does LightGBM handle the UNCTAD NaN natively, and what does that mean?
  \item How does the what-if chat re-lookup the UNCTAD benchmark when the lane changes?
  \item What's the RMSE in USD terms for a typical shipment?
  \item What happens when the user asks something outside the corpus?
  \item What are the limitations of using health commodity data for general freight pricing?
  \item How would you retrain this model on new data in production?
  \item Why is random train/test split wrong for this data?
  \item What would you add with 6 more months?
  \item What's the coverage rate of UNCTAD on USAID lanes, and what's your fallback strategy?
\end{itemize}

% ── 5. Alternatives Table ────────────────────────────────────────────────────
\section{Alternatives Table}

Debug max 2 hours before swapping.

\begin{tabularx}{\textwidth}{lXX}
\toprule
\textbf{Blocker} & \textbf{Primary} & \textbf{Alternative} \\
\midrule
UNCTAD coverage too low ($<$30\% of USAID rows) & UNCTAD join & Drop UNCTAD feature. USAID-only model. Reframe README: ``pure health commodity SC analysis.'' \\
\addlinespace
UNCTAD rates uncorrelated with USAID costs & Use as feature & Drop UNCTAD feature. Keep model simple. Still a strong portfolio piece. \\
\addlinespace
USAID too dirty/small & USAID Kaggle & DataCo Smart SC (note: shipping modes are fulfillment tiers, not real freight modes --- reframe entirely if used) \\
\addlinespace
Groq rate limits hit during demo & Groq free & together.ai free tier or OpenAI \$5 budget \\
\addlinespace
FAISS install fails (Windows) & faiss-cpu & chromadb --- cleaner pip, persistent by default \\
\addlinespace
SHAP slow on LightGBM & \texttt{shap.TreeExplainer} & Gain-based feature importance from LightGBM directly (less accurate but fast) \\
\addlinespace
What-if intent parsing unreliable & Groq LLM parse & Rule-based keyword fallback for common cases (mode\_change, route\_change) \\
\addlinespace
HF Spaces too slow & HF Spaces & Streamlit Community Cloud \\
\addlinespace
PDF parsing chaotic & pypdf & Convert corpus to .txt manually, commit, parse as plain text \\
\addlinespace
sentence-transformers slow & all-MiniLM-L6-v2 & bge-small-en-v1.5 (smaller, faster on CPU) \\
\addlinespace
16GB RAM pressure & float32 embeddings & Cast to float16; FAISS IVF index \\
\bottomrule
\end{tabularx}

% ── 6. Risk Flags ────────────────────────────────────────────────────────────
\section{Risk Flags}

\begin{warnbox}
\textbf{1. Time-based split is mandatory.} USAID has a date column. Use time-based split (80/20 at Dec 2013). Random split on temporal data = silent leakage. Will come up in interviews.
\end{warnbox}

\begin{warnbox}
\textbf{2. SHAP grounds the explanation --- without it, LLM hallucinates.} Do not skip SHAP. An LLM explaining a prediction without being shown the actual model drivers is guessing. This is the most common mistake in ``explainable AI'' demos and interviewers notice.
\end{warnbox}

\begin{warnbox}
\textbf{3. UNCTAD temporal mismatch must be disclosed.} USAID: 2006--2015. UNCTAD: 2016--2021. No year overlap. UNCTAD is a structural benchmark, not a time-matched join. State this explicitly in README and interviews. The defensible framing: bilateral freight rate intensities are structural signals that change slowly; using the 2016--2021 median as a lane benchmark is standard practice in freight gravity-model research.
\end{warnbox}

\begin{warnbox}
\textbf{4. UNCTAD coverage audit is Week 2's most important output.} If coverage is thin, the UNCTAD feature may hurt more than it helps (high NaN rate + weak signal = noisy feature). Run the audit before building the full pipeline. Know your fallback hierarchy: HS4 match $\rightarrow$ HS2 aggregate $\rightarrow$ mode-level country-pair average $\rightarrow$ NaN (LightGBM handles).
\end{warnbox}

\begin{warnbox}
\textbf{5. What-if intent parsing is the hardest piece.} LLM-based intent extraction can be flaky. Test with 20+ diverse messages. Build the rule-based fallback from the alternatives table. Two LLM calls per what-if query = double rate limit exposure. Cache aggressively.
\end{warnbox}

\begin{warnbox}
\textbf{6. st.session\_state in Streamlit.} Stateful chat requires careful session state management. Streamlit reruns the entire script on every interaction. The shipment dict and chat history must live in \texttt{st.session\_state} or they reset. Test before Week 10 build-out.
\end{warnbox}

\begin{warnbox}
\textbf{7. Week 6 buffer is non-negotiable.} Do not skip it when feeling ahead in Week 4.
\end{warnbox}

\begin{warnbox}
\textbf{8. Groq TOS.} Verify Groq free tier allows public portfolio demo use before launch.
\end{warnbox}

% ── 7. What NOT To Do ────────────────────────────────────────────────────────
\section{What NOT To Do}

\begin{itemize}[noitemsep]
  \item \textbf{No deep learning on tabular data.} Red flag in SC interviews.
  \item \textbf{No Streamlit before Week 10.} Model and explanation must be solid first.
  \item \textbf{No Docker.} HF Spaces handles it.
  \item \textbf{No MLflow/W\&B.} Mention in README next steps, do not implement.
  \item \textbf{No custom embedding fine-tuning.} Out of scope.
  \item \textbf{No random split if data is temporal.} See Risk Flag 1.
  \item \textbf{Do not let the LLM explain without SHAP input.} See Risk Flag 2.
  \item \textbf{Do not pretend UNCTAD is a time-matched join.} See Risk Flag 3.
  \item \textbf{Do not force UNCTAD into the model if coverage audit fails.} See Risk Flag 4 and Alternatives Table.
\end{itemize}

% ── 8. Success Criteria ──────────────────────────────────────────────────────
\section{Success Criteria}

Project A1 is done when all of the following are true:

\begin{enumerate}
  \item Deployed Streamlit app on HF Spaces, live and accessible via public URL.
  \item Tab 1: predicts cost from shipment form; UNCTAD benchmark shown alongside prediction.
  \item Tab 1: plain-English explanation with SHAP feature table, UNCTAD callout, and cited trade doc context.
  \item Tab 2: chat understands current shipment context from session state.
  \item Tab 2: what-if queries rerun model (with updated UNCTAD lookup if lane changed) and return new prediction + delta + explanation.
  \item Tab 2: general trade questions answered via RAG with citations.
  \item GitHub repo has clean README with results table (RMSE/MAE baseline vs LightGBM), UNCTAD coverage stats, demo link, architecture description, two screenshots.
  \item You can answer 14 interview questions about this project without notes, including the UNCTAD temporal mismatch question.
  \item LightGBM beats linear baseline on RMSE.
  \item OR: UNCTAD feature dropped after coverage audit fails --- USAID-only model shipped, README updated honestly.
\end{enumerate}

\begin{infobox}
\textbf{Next checkpoint:} Week 2 --- post UNCTAD coverage audit results: what \% of USAID rows get direct match, what \% fallback, what \% NaN. Post correlation between UNCTAD rates and USAID actual costs. Go/no-go decision on UNCTAD feature before building pipeline.
\end{infobox}

\end{document}                                                                                                                                                                                                                                                                                                                                                                                                                                                                                                                                                                                                                                                                                                                                                                                                                                                                                                                                                                                                                                                                                                                                                                                                                                                                                                                                                                                                                                                                                                                                                                                                                                                                                                                                                                                                                                                                                                                                                                                                                                                                                                                                                                                                                                                                                                                                                                                                                                                                                                                                                                                                                                                                                                                                                                                                                                                                                                                                                                                                                                                                                                                                                                                                                                                                                                              ```````````````````````````````````````````````````````````````````````````````````````````````````````````````````````````````````````````````````````````````````````````````````````````````````````````````````````````````````````````````````````````````````````````````````````````````````````````````````````````````````````````````````````````````````````````````````````````````````````````````````````````````````````````````````````````````````````````````````````````````````````````````````````````````````````````````````````````````````````````````````````````````````````````T6'''''''''''''''''''''''''''''''''''''''''''''''''''''''''''''''''''''''''''''''''''''''''''''''''''''''''''''''''''''''''''''''''''''''''''''''''''